\section{One polynomial block-cost criterion without proper boundary truncation}
\label{sec:polynomial-partial-return}

The all-compact-test criterion admits a further simplification
when the Hilbert quotient retains or discards each primary
factor in full.  In this case the finite return has no
algebraic freedom.  A single polynomial estimate for its
actual norm is equivalent to returns for every smooth test.
The estimate itself is not established here at a general width.

Retain all fixed data and the original adapted block norms
of MA-PR-006.  Assume explicitly that
\[
                 r_\rho\in\{0,m_\rho\}
                    \qquad\hbox{for every zero }\rho.
 \tag{MA-PR.100}
\]
This holds automatically when \(s\le1\), because then
no edge jet is continuous, and at widths with no boundary
zero for every \(s\ge0\).  It also permits fully retained
boundary factors.  It is not automatic at other widths
with \(s>1\).  Under (MA-PR.100), eligibility imposes
no condition: \(\mathcal E_{\rm el}=\mathcal E\).

The primary decomposition supplies the canonical module
inclusion \(\iota_n:B_n\to J_n\): it is the identity
on every fully retained factor and zero at every
discarded factor.  Define its measured whole-quotient cost by
\[
       \mathcal C_n
           =\|\iota_n\mathscr R_n\|_{\mathcal L(Q_H,J_n)}.
 \tag{MA-PR.101}
\]
Set \(\mathcal C_n=0\) if \(B_n=0\).
The target has its original holomorphic quotient norm,
and \(\mathscr R_n\) acts on the whole \(Q_H\).
Thus (MA-PR.101) includes every oblique-projection cost.

\paragraph{MA-PR-010 (a single polynomial gate for all tests).}
Under (MA-PR.100), the following conditions are equivalent.
\begin{enumerate}
\item There are fixed finite constants \(C>0\), \(D\ge0\) with
\[
                \mathcal C_n\le C\nu_n^D
                           \qquad(n\in\mathbb Z).
 \tag{MA-PR.102}
\]
\item Every \(g\in\mathcal E\) has a continuous
equivariant return satisfying (MA-PR.51).
\item Every \(g\in C_c^\infty(\mathbb R)\) with
\(g(0)=0\) has such a return.
\item Every \(g\in\mathcal D_{[1,2]}\) has such a return.
\end{enumerate}
The same equivalence holds if condition~3 is stated
for all compact smooth tests without the value-zero
restriction.  Condition~4 concerns the entire fixed-support
Fr\'echet space, not just a chosen sequence of tests.

Whenever these conditions hold, the return \(S_g\) is
unique as a map to \(Q_\infty\), is linear in \(g\),
and has an \(\mathcal E\)-valued lift such that
\[
       (g,v)\longmapsto\widetilde S_gv
             \quad\hbox{is jointly continuous on }
                         \mathcal E\times Q_H .
 \tag{MA-PR.103}
\]
Its common nuclear representations have coefficients
in every positive \(\ell^p\), uniformly for \(g\)
in bounded subsets of \(\mathcal E\).
The quotient-valued return is unique; its representative
lift into \(\mathcal E\) uses the chosen MT-M18 lifts.
All constants concern the fixed function, Hilbert
parameters and adapted blocks.  No uniformity as these
data vary is asserted.

If (MA-PR.102) fails, then
\[
 \{g\in\mathcal D_{[1,2]}:
                   \sup_n c_n(g)=\infty\}
       \quad\hbox{is a dense }G_\delta
                    \hbox{ in }\mathcal D_{[1,2]}.
 \tag{MA-PR.104}
\]
Every such compact test satisfies \(g(0)=0\) and
has no continuous equivariant return.  This alternative
does not assert the arithmetic failure of (MA-PR.102).

\begin{proof}
\emph{Uniqueness of the finite map and exact multiplication.}
At a fully retained factor the reduction is the identity
\(\mathbb C[z]/(z^m)\to\mathbb C[z]/(z^m)\).
A module map factoring convolution is therefore forced
to be multiplication by the full Taylor class.
At a discarded factor there is no source factor;
all maps from other primary supports to it vanish by
coprimality of their annihilating polynomials.
Thus, for any \(b\in J_n\), the sole feasible map
with reduction \(M_{\pi_nb}\) is
\[
       L_n(b)=M_b\iota_n=\iota_nM_{\pi_nb}.
 \tag{MA-PR.105}
\]
The equality follows factor by factor in the actual
finite modules; it does not use an orthogonal Hilbert
decomposition.  In particular
\[
 \begin{aligned}
 c_n(g)&=\|M_{b_n(g)}\iota_n\mathscr R_n\|,\\
 c_n(g)&\le\|b_n(g)\|_n\mathcal C_n.
 \end{aligned}
 \tag{MA-PR.106}
\]
The latter estimate is the submultiplicative original
quotient norm on \(J_n\).
Every continuous global return has exactly these
block maps by (MA-PR.57).  Since the full smooth
jet kernel is zero in \(Q_\infty\), they determine
the global return uniquely whenever it exists.

\emph{The polynomial estimate gives all smooth returns.}
Assume (MA-PR.102).  For every \(g\in\mathcal E\),
the sequence \(\|b_n(g)\|_n\) decreases faster than
every power.  Equation (MA-PR.106) therefore satisfies
the exact return criterion (MA-PR.52).
More explicitly, choose an integer \(N_0>3+|a|\).
Integration by parts on the fixed strip containing
all full block contours gives, for each integer \(M\ge0\),
\[
 \|b_n(g)\|_n\le C_{a,M}\nu_n^{-M}
                  \sum_{l=0}^{M}p_{N_0,l}(g).
 \tag{MA-PR.107}
\]
At large \(|n|\), apply \(M\) integrations by parts
to \(e^{(\operatorname{Re}w-a)x}g(x)\) in its Fourier
integral.  Its differentiated integrands are bounded
by the displayed exponential seminorms times an
integrable exponential in \(|x|\).
The finitely many small heights follow from the
undifferentiated integral after enlarging the constant.
This proves (MA-PR.107) for every derivative order;
it is an estimate in the original quotient norm.

Set
\[
 \begin{aligned}
 F_{n,g}&=M_{b_n(g)}\iota_n\mathscr R_n,\\
 \widetilde S_gv&=\sum_n\mathcal T_nF_{n,g}v,
             \qquad S_g=q_\infty\widetilde S_g .
 \end{aligned}
 \tag{MA-PR.108}
\]
The finite map and \(\mathcal T_n\) are fixed,
so the series is linear in \(g\) and in \(v\).
For a defining target seminorm \(p_{N,j}\), let
\(D_{N,j}\) be its polynomial lift exponent in
(MA-PR.44).  Choose an integer
\(M>D+D_{N,j}+2\).
Then (MA-PR.102), (MA-PR.107) and the lift bound give
\[
 p_{N,j}(\widetilde S_gv)
       \le C_{N,j,M}\left(\sum_{l=0}^{M}
                                   p_{N_0,l}(g)\right)\|v\|.
 \tag{MA-PR.109}
\]
The series converges absolutely in every seminorm,
uniformly on the corresponding bounded product sets.
This proves the claimed joint continuity and the
actual lift.  Its full block class is \(F_{n,g}v\).
Taking retained jets proves \(iS_g=C_g^{Q_H}\) by
the actual common-kernel theorem MA-WP-001.
The module identities in (MA-PR.105) prove the
centered equivariance in (MA-PR.51).

For a bounded test set \(\mathcal B\subset\mathcal E\),
the sequence
\(\beta_n=e^{-\nu_n}+\sup_{g\in\mathcal B}\|F_{n,g}\|\)
is positive and rapidly decreasing by (MA-PR.107).
Using an Auerbach basis of \(J_n\), take target
vectors \(\beta_n^{1/2}\mathcal T_ne_{n,l}\),
coefficients \(\beta_n^{1/2}\), and functionals
\(\ell_{n,l}F_{n,g}/\beta_n\).
The target vectors form one bounded family in
\(\mathcal E\), the source norms are at most one,
and the coefficients repeated \(d_n=O(\log\nu_n)\)
times belong to every positive \(\ell^p\).
This proves the uniform all-order nuclear assertion.
Condition~1 has now given condition~2 and the
stated coherent topology.  The implications
\(2\Rightarrow3\Rightarrow4\) follow from the
test-space inclusions.

\emph{One fixed positive bump with uniform inverse classes.}
For the converse fix \(x_0=3/2\) and the compact
rectangle
\[
 \mathcal R_a=[-1-a,2-a]
                           +i[1/4,7/4].
 \tag{MA-PR.110}
\]
Let \(M_a=\max_{z\in\mathcal R_a}|z|\).
Choose \(0<\epsilon<1/4\) so small that
\(e^{M_a\epsilon}-1\le1/2\).
Take a nonnegative \(\phi\in C_c^\infty
((x_0-\epsilon,x_0+\epsilon))\) with integral one.
For \(z\in\mathcal R_a\), the elementary exponential
estimate gives
\[
 \begin{split}
 \left|e^{-x_0z}\mathcal M\phi(z)-1\right|
   &=\left|\int\bigl(e^{z(x-x_0)}-1\bigr)\phi(x)\,dx\right|\\
   &\le e^{M_a\epsilon}-1\le\frac12 .
 \end{split}
 \tag{MA-PR.111}
\]
In particular the transform is bounded above and
away from zero on this whole rectangle, with
constants depending only on the fixed \(a\) and bump.

Put \(g_n(x)=e^{-inx}\phi(x)\).
It is supported in the same compact subset of
\((1,2)\) and has \(g_n(0)=0\).
For \(w\in\overline\Omega_n\), the selected heights
give \(w-a-in\in\mathcal R_a\), and
\[
 \begin{aligned}
 b_n(g_n)(w)&=\mathcal M\phi(w-a-in),\\
 \|b_n(g_n)\|_n+\|b_n(g_n)^{-1}\|_n&\le C_a .
 \end{aligned}
 \tag{MA-PR.112}
\]
For the inverse, use the actual holomorphic function
\(1/\mathcal M\phi(w-a-in)\) on a neighborhood of
the closed rectangle, and take its class in \(J_n\).
Equation (MA-PR.111) makes its supremum norm
uniform in \(n\).  There is no inverse Hermite
matrix or zero-separation factor in this argument.
Blocks with no retained factor require no inverse
estimate, since their cost is zero.

Every derivative of the modulated bump has the
explicit formula
\[
 \begin{aligned}
 g_n^{(l)}(x)
   &=e^{-inx}\sum_{q=0}^{l}\binom lq
                             (-in)^{l-q}\phi^{(q)}(x),\\
 p_j^{[1,2]}(g_n)&\le C_{\phi,j}\nu_n^j,
       \qquad
 p_{N,j}(g_n)\le C_{\phi,N,j}\nu_n^j .
 \end{aligned}
 \tag{MA-PR.113}
\]
The last bound also follows directly because
the common support is contained in \([1,2]\).
Constants may depend on the displayed derivative
orders, but never on \(n\).

\emph{All fixed-support returns force the polynomial cost.}
Assume condition~4.  There are no eligibility
constraints under (MA-PR.100), so MA-PR-009 with
\(X_K=\mathcal D_{[1,2]}\) and \(k=0\) gives
one finite \(j\) and one \(C\) such that
\[
                     \sup_m c_m(g)\le C p_j^{[1,2]}(g)
                          \qquad(g\in\mathcal D_{[1,2]}).
 \tag{MA-PR.114}
\]
This uniform conclusion uses the return hypothesis
for the entire Fr\'echet space.  Existence merely
for the sequence \(g_n\) would not justify it.

Use the unique finite map (MA-PR.105) and multiply
it by the inverse full class in (MA-PR.112).
Submultiplicativity and (MA-PR.113)--(MA-PR.114) give
\[
 \begin{aligned}
 \mathcal C_n
   &=\|M_{b_n(g_n)^{-1}}F_{n,g_n}\|\\
   &\le C_a c_n(g_n)
    \le C_a C p_j^{[1,2]}(g_n)
    \le C'\nu_n^j .
 \end{aligned}
 \tag{MA-PR.115}
\]
Thus condition~1 holds.  This completes all four
equivalences.  If condition~3 omits the value-zero
restriction, condition~2 still implies it, and it
still implies condition~4, so that version is
equivalent as well.

If the polynomial estimate fails, (MA-PR.115)
shows that no bound of the form (MA-PR.114) can
hold for any finite \(j,C\).
The \(k=0\) failure alternative of MA-PR-009
then gives precisely (MA-PR.104).
No arithmetic instance of that failure is needed
for this conditional implication.
\end{proof}

\paragraph{Trace transport and the remaining domains.}
Under these equivalent conditions, the common nuclear
series in (MA-PR.108) permits the actual trace passage
\[
 \begin{gathered}
 B_g=S_g i,\qquad E_g=C_{T_{-a}g}^{Q_\infty}-B_g,\\
 \begin{aligned}
 \operatorname{Tr}_{Q_H}C_g
   &=\operatorname{Tr}_{Q_\infty}B_g,\\
 \operatorname{Tr}_{Q_\infty}B_g
   &=\operatorname{Tr}_{Q_\infty}C_{T_{-a}g}
                         -\operatorname{Tr}_{Q_\infty}E_g .
 \end{aligned}
 \end{gathered}
 \tag{MA-PR.116}
\]
The rank-one pairings coincide and are absolutely
summable, exactly as proved in MA-PR.72.
In each full retained factor \(B_g\) is multiplication
by the full class; in each discarded factor it is zero.
Thus \(E_g\) is an actual nuclear correction on the
discarded smooth factors.  For an original compact
test with \(g(0)=0\) in the rational centered sector, the existing
Meyer identification and MA-WP-003 supply the same
geometric transport as MA-PR.62, now with
\(\operatorname{Tr}E_g\) itself instead of a
regularized correction limit.
The arithmetic disappearance of that correction
has not been proved.

These return maps act on the specified smooth and
Hilbert quotients.  They do not extend the original
off-diagonal maps of Meyer's common-space projection
to new Hilbert domains, or identify that projection
with a source PW cutoff.  No ambient trace-class
claim for arbitrary noncompact \(g\in\mathcal E\)
is made by (MA-PR.116).
The assumption \(s\le1\) removes proper edge
truncation; it does not by itself prove
(MA-PR.102).  Likewise the theorem does not supply
the polynomial cost bound merely from absence of
an exact boundary zero.

The exact inputs are the actual criterion and
whole-quotient norms MA-PR-006, the proved Baire
implication MA-PR-009, and the full adapted
MT-M18 block topology and entire-lift estimates.
The modulated-bump inverse estimate is proved here
in the original holomorphic norm.  No A, M,
simplicity, generic-edge distance estimate, or
new outside preprint is used.

The finite-length \(\mathbb C[w]\)-module and nuclear-trace
placement is the one specified at (MA-PR.80) and in
MA-PR-009.  Under (MA-PR.100), the primary decomposition
itself supplies the finite retained/discarded inclusion;
its global analytic cost is still the separate estimate
(MA-PR.102).  A Grothendieck-group relation alone supplies
no such estimate.  The inherited cote~1 scan locators
and exterior-power determinant conventions retain their
previously stated scope; no new reading of those scans
or arithmetic positivity conclusion is claimed.
